\section*{Author Section 3.1: setup}
Let $p$ be the hard distribution for $Q^1_{\epsilon+12\zeta}(f)$ from Yao's lemma, so $Q^1_{\epsilon+12\zeta}(f)=Q^1_{p,\epsilon+12\zeta}(f)$. Extend $f$ to $\widetilde f\subseteq\mathcal X\times(\mathcal Y\cup\{y^*\})\times\mathcal Z$ by including $(x,y^*,z)$ for every $x,z$. Extend $p$ with $p(y^*)=0$. Then
\[Q^1_{p,\gamma}(\widetilde f)=Q^1_{p,\gamma}(f)\tag{1}\]
for every error $\gamma$. Also
\[Q^1_\gamma(f^k)\ge Q^1_\gamma(\widetilde f^k).\tag{2}\]
Every protocol for $f^k$ gives one for $\widetilde f^k$: on positions with $y^*$ Bob pretends an input from $\mathcal Y$, runs it and answers accordingly. Its output is correct there regardless. We show, for the distribution $q$ below,
\[Q^1_{q^k,1-(1-\epsilon)^{\Omega(\zeta^6k/\log|\mathcal Z|)}}(\widetilde f^k)\ge\frac{\zeta^5k}{60}Q^1_{p,\epsilon+12\zeta}(\widetilde f)-k\log\log\frac{24}{5\zeta}.\tag{3}\]
Since $Q^1_\gamma(\widetilde f^k)\ge Q^1_{q^k,\gamma}(\widetilde f^k)$, (1)--(3) imply Theorem1. Set $q(x,y)=(1-\zeta)p(x,y)$ for $y\in\mathcal Y$ and $q(x,y^*)=\zeta p(x)$. Clearly $q(x,y^*)=q(x)q(y^*)$ and
\[\|p(x,y)-q(x,y)\|_1\le2\zeta.\tag{4}\]
For each $i\in[k]$, define independent copies of $X_iY_iD_iG_i$, with marginal $X_iY_i$ equal to $q$. $D_iG_i$ are correlation-breaking variables: given their value, $X_i,Y_i$ are independent. $D_i$ is uniform in $\{0,1\}$. Its conditional $G_i$ distribution is:
\begin{itemize}\item If $D_i=0$, $G_i=x$ w.p. $p(x)$.
\item If $D_i=1$, $G_i=y^*$ w.p. $1-(1-\zeta)^{2/3}$, and $G_i=y$ w.p. $(1-\zeta)^{2/3}p(y)$.\end{itemize}
Conditional $X_iY_i$ distributions are:
\begin{itemize}\item If $D_i=0,G_i=x$: $(x,y^*)$ has probability $\zeta$, and $(x,y)$ has probability $(1-\zeta)p(y|x)$.
\item If $D_i=1,G_i=y^*$: $(x,y^*)$ has probability $p(x)$.
\item If $D_i=1,G_i=y$: $(x,y^*)$ has probability $(1-(1-\zeta)^{1/3})p(x|y)$, and $(x,y)$ has probability $(1-\zeta)^{1/3}p(x|y)$.\end{itemize}
\paragraph{Lemma 30.} $P_{X_iY_i}(x,y)=q(x,y)$ for every $(x,y)\in\mathcal X\times(\mathcal Y\cup\{y^*\})$.
\begin{proof}For $D_i=0$, $G_i=X_i$, so $P_{G_iY_i|D_i=0}=P_{X_iY_i|D_i=0}=q$. For $D_i=1$ and $y\in\mathcal Y$, $Y_i=y$ only if $G_i=y$, hence
\begin{align*}P_{X_iY_i|D_i=1}(x,y)&=P_{G_i|D_i=1}(y)P_{X_iY_i|D_i=1,G_i=y}(x,y)\\&=(1-\zeta)^{2/3}p(y)(1-\zeta)^{1/3}p(x|y)=(1-\zeta)p(x,y)=q(x,y).\end{align*}
$Y_i=y^*$ can arise when $G_i=y^*$ or any $y\in\mathcal Y$, so
\begin{align*}P_{X_iY_i|D_i=1}(x,y^*)&=P_{G_i|D_i=1}(y^*)P_{X_iY_i|D_i=1,G_i=y^*}(x,y^*)\\&\quad+\sum_yP_{G_i|D_i=1}(y)P_{X_iY_i|D_i=1,G_i=y}(x,y^*)\\&=(1-(1-\zeta)^{2/3})p(x)\\&\quad+(1-\zeta)^{2/3}(1-(1-\zeta)^{1/3})\sum_yp(y)p(x|y)\\&=(1-(1-\zeta)^{2/3})p(x)+((1-\zeta)^{2/3}-(1-\zeta))p(x)\\&=\zeta p(x)=q(x,y^*).\end{align*}\end{proof}
In particular $P_{X_iY_i}(x,y^*)=P_{X_i}(x)P_{Y_i}(y^*)$. Also
\[P_{Y_iG_i|D_i=1}(Y_i\ne G_i)=(1-\zeta)^{2/3}(1-(1-\zeta)^{1/3})\le1-2\zeta/3-1+\zeta=\zeta/3.\tag{5}\]
Let $\mathcal P$ be a one-way protocol for $\widetilde f^k$ of cost $ck$. Alice/Bob's inputs are $X=X_1\cdots X_k$, $Y=Y_1\cdots Y_k$. They share an input-uncorrelated entangled pure state on $E^AE^B$. Alice applies $V^A$ on $XE^A$, producing message $M$ and discarded register $A$. Let $\ket\theta_{AME^B|x}$ denote that pure state; it only depends on $y$ via $x$. For input distribution $P_{XY}$ the CQ state is
\[\theta_{XYAME^B}=\sum_{xy}P_{XY}(xy)\proj{xy}_{XY}\otimes\proj\theta_{AME^B|x}.\]
Purify it using $\widetilde X,\widetilde Y$:
\[\ket\theta_{X\widetilde XY\widetilde YAME^B}=\sum_{xy}\sqrt{P_{XY}(xy)}\ket{xxyy}_{X\widetilde XY\widetilde Y}\ket\theta_{AME^B|x}.\]
Bob applies $V^B$ on $YME^B$, converting $ME^B$ to $BZ$, with output $Z=Z_1\cdots Z_k$. Call the resulting pure state $\ket\rho_{X\widetilde XY\widetilde YABZ}$. Define $P_{XYDGZ}$ with the earlier $XYDG$, $Z$ independent of $DG$ given $XY$, and its conditional law from measuring $Z$ in $\ket\rho$.
\begin{center}\includegraphics[width=0.72\linewidth]{figure1.png}\\Figure1. One-way quantum protocol $\mathcal P$.\end{center}
